\begin{enumerate}

    \item 如本节所述的柯尼斯堡漫步，或者更一般地，在一个任意图中，精确地穿过每条边一次并起止于同一节点的回路，被称为\index{Eulerian circuit}\emph{欧拉回路}。而\index{Eulerian path}\emph{欧拉路径}也精确地穿过图的每条边一次，但它的起点和终点是不同的节点。对于下面的每一个图，判断是否存在欧拉回路或欧拉路径，如果存在，请画出它。
    
    \begin{center}
    \input{figures/Euler_circuit_problems_a.tex}
    \end{center}
    
    \begin{center}
    \input{figures/Euler_circuit_problems_b.tex}
    \end{center}
    
    \item 完成“每个图都有偶数个奇数度节点”这一事实的证明。
    
    \wbvfill
    
    \item 请给出一个论证，说明为什么一个$8 \times 8$的棋盘，在对角线两端各去掉一个方格后，无法用多米诺骨牌完全覆盖。
    
    \wbvfill
    
    \workbookpagebreak
    
    \item 证明，如果$n$是奇数，任何$n \times n$的棋盘，在去掉一个与其角格颜色相同的方格后，可以用多米诺骨牌完全覆盖。
    \wbvfill
    
    \item 五个\index{tetromino}四格骨牌（俄罗斯方块玩家很熟悉）是多米诺骨牌的亲戚，由四个小方格组成。
    \begin{center}
    \input{figures/five_tetrominoes.tex}
    \end{center}
    
    \noindent 这五个四格骨牌总共包含20个方格，所以可以想象它们可以用来铺满一个$4 \times 5$的棋盘。证明这实际上是不可能的。
    
    \wbvfill
    
    \workbookpagebreak
    
    \item 陈述图中存在欧拉回路的充要条件。
    \wbvfill
    
    \item  陈述图中存在欧拉路径的充要条件。
    \wbvfill
    
    \newpage
    
    \item 构建4阶和5阶的幻方。
    
    \wbvfill
    
    \workbookpagebreak
    
    \item 一个2阶的幻六边形需要将数字1到7填入下面的六边形阵列中。幻方的条件意味着每9条相邻六边形的“线”上的数字之和都相同。这可能吗？
    \begin{center}
    \input{figures/magic_hexagon.tex}
    \end{center}
    
    \wbvfill
    
    \item 存在3阶的幻六边形吗？
    
    \wbvfill
    
    \end{enumerate}
    
    
    
    %% Emacs customization
    %% 
    %% Local Variables: ***
    %% TeX-master: "GIAM-hw.tex" ***
    %% comment-column:0 ***
    %% comment-start: "%% "  ***
    %% comment-end:"***" ***
    %% End: ***